\section{Microphone Arrays: Localization of Wide-Band Sources}

\subsection{Introduction}
Let us focus on a scenario involving a single acoustic source, small or moderate reverberation, and a \textbf{wide-band source signal}.
In this scenario, spatial filtering or parametric methods may not represent the most efficient solution.
We will study two techniques for source localization: the \textit{Generalized Cross-Correlation} method and the \textit{Adaptive Eigenvalue Decomposition} method.

\subsection{Data model}
When we consider wide-band sources, we cannot rely on the sinusoidal model.
For simplicity, without loss of generality, we consider a pair of microphones.
Let us define the \textbf{source} signal $s$, the \textbf{channel} response $g_k$ from source to $k$-th microphone and the microphone \textbf{additive noise} $e_k$.
The wide-band array model describes the microphone signal $y_k$ as:

\[
y_k(n) = g_k * s \ (n) + e_k(n),
\qquad k = 1, 2
\]

We may model the two channel (impulse) responses as follows:

\[
g_1(n) = \underbrace{\frac{1}{r_1} \delta(n - \tau_1)}_{\text{direct signal}} + \underbrace{\sum_i \alpha_i \ \delta(n - \tau_{1i})}_{\text{early reflections}},
\qquad
g_2(n) = \underbrace{\frac{1}{r_2} \delta(n - \tau_2)}_{\text{direct signal}} + \underbrace{\sum_i \beta_i \ \delta(n - \tau_{2i})}_{\text{early reflections}}
\]

The information about the source position is contained in the propagation delays $\tau_1$ and $\tau_2$.

\begin{figure}[H]
    \centering
    \includegraphics[width=0.3\linewidth]{Two_microphones DOA.png}
    \caption{DOA in a two-microphones configuration with wide-band source in far-field}
\end{figure}

The source emission time is generally unknown, so we aim at estimating the \textbf{Time Difference Of Arrival (TDOA)}, also known as the time delay:

\[
\tau_{12} = \tau_2 - \tau_1
\]

Assuming that the source is in the far-field with respect to the microphone pair, the TDOA is related to the DOA through a simple trigonometric relationship:

\[
\tau_{12} = \frac{d \sin \left( \theta \right)}{c}
\]

Given an estimate $\hat{\tau}_{12}$ of the TDOA, the estimate $\hat{\theta}$ of the DOA is computed as:

\[
\hat{\theta} = \arcsin \left( \frac{\hat{\tau}_{12} \ c}{d} \right)
\]


The problem of DOA estimation can be brought to the \textbf{problem of estimating the TDOA} between the microphone pair.
Different approaches can be used to estimate TDOAs.
The first approach finds the maxima of the \textbf{cross-correlation} of the microphone signals, while the second approach identifies the \textbf{channel responses} $g_{1}$ and $g_{2}$.

\subsection{TDOA estimation: Generalized Cross-Correlation (GCC)}
We expect the two microphones to record two delayed copies of the source signal, along with some disturbs (i.e. environmental noise and reverberation).
If the amount of noise and reverberation is low compared to the strength of the source signal, we can cross-correlate the microphone signals to estimate the TDOA.
The \textbf{cross-correlation} of the microphone signals is given by the inverse Fourier Transform of the cross-spectrum:

\[
R_{12}(t) = \frac{1}{2\pi} \int_{-\infty}^{\infty} X_{1}(\omega) X_{2}^{*}(\omega) \ e^{j\omega t} \ d\omega
\]

The TDOA estimation $\hat{\tau}_{12}$ is the time-lag that \textbf{maximizes the cross-correlation} function, within a valid time interval:

\[
\hat{\tau}_{12} =\underset{\tau \in \mathcal{D}}{\mathrm{arg \ max}} \left( R_{12}(\tau) \right), 
\qquad 
\mathcal{D} = \left[ -\frac{d}{c}, \frac{d}{c} \right]
\]

The \textbf{problem} is that, if the signal presents some \textbf{periodicities} (i.e. it is somehow harmonic), the cross-correlation function will present peaks that arise from the auto-correlation of the signal.
These peaks may overcome the searched one, leading to a wrong TDOA estimation.

\begin{figure}[H]
    \centering
    \includegraphics[width=0.5\linewidth]{Cross-correlation and GCC.png}
    \caption{TDOA estimation: comparison between cross-correlation (blue) and generalized cross-correlation (orange)}
\end{figure}

Due to the impulsive behaviour of the auto-correlation function, the overlap is minimized for white signals.
The signal spectrum can be \textbf{artificially whitened} by introducing a weight function in the cross-correlation function.
We talk about \textbf{Generalized Cross-Correlation (GCC)}:

\[
R_{12}(t) = \frac{1}{2\pi} \int_{-\infty}^{\infty} \psi_{12}(\omega) X_{1}(\omega) X_{2}^{*}(\omega) \ e^{j\omega t} \ d\omega
\]

\[
\psi_{12}(\omega) = \frac{1}{\abs{X_{1}(\omega) X_{2}^{*}(\omega)}}: \ \text{PHAT (Phase Transform) function}
\]

\clearpage

\subsection{TDOA estimation: Adaptive Eigenvalue Decomposition (AED)}
An alternative approach for the TDOA estimation starts from the identification of the channel responses $g_k$.
The first peak in the impulse responses is related to the \textbf{direct path} from the source to the microphones, while we exclude the reverberant components (i.e. early reflections and reverberation tail).
We recover the TDOA as the difference between the time-lags related to the direct paths at the two microphones.

\begin{figure}[H]
    \centering
    \includegraphics[width=0.5\linewidth]{AED impulse response.png}
    \caption{Example of channel response $g_k(n)$}
\end{figure}

\textbf{Adaptive Eigenvalue Decomposition (AED)} is based on the following result, valid in a noise-free scenario:

\[
y_{1}(n) * g_{2}(n) = s(n) * g_{1}(n) * g_{2}(n) = y_{2}(n) * g_{1}(n)
\]

This expression can be approximated considering $M$-long frames of the microphone signal:

\[
\underline{\mathrm{y}}_{k}(n) =
\begin{bmatrix}
y_{k}(n) \\
y_{k}(n-1) \\
\vdots \\
y_{k}(n-M+1)
\end{bmatrix},
\qquad
\underline{\mathrm{g}}_{k} =
\begin{bmatrix}
g_{k}(0) \\
g_{k}(1) \\
\vdots \\
g_{k}(M-1)
\end{bmatrix},
\qquad
k=1,2
\]

\[
\underline{\mathrm{y}}_{1}^{T}(n) \ \underline{\mathrm{g}}_{2} = \underline{\mathrm{y}}_{2}^{T}(n) \ \underline{\mathrm{g}}_{1}
\quad \Rightarrow \quad
\underline{\mathrm{y}}_{1}^{T}(n) \ \underline{\mathrm{g}}_{2} - \underline{\mathrm{y}}_{2}^{T}(n) \ \underline{\mathrm{g}}_{1} = 0
\]

We can rewrite the previous expression in a statistical fashion.
We define the following:

\[
\left[ \mathrm{R} \right] =
\begin{bmatrix}
E\left\{ \underline{\mathrm{y}}_{1}(n) \ \underline{\mathrm{y}}_{1}^{T}(n) \right\} &
E\left\{ \underline{\mathrm{y}}_{1}(n) \ \underline{\mathrm{y}}_{2}^{T}(n) \right\} \\[8pt]
E\left\{ \underline{\mathrm{y}}_{2}(n) \ \underline{\mathrm{y}}_{1}^{T}(n) \right\} &
E\left\{ \underline{\mathrm{y}}_{2}(n) \ \underline{\mathrm{y}}_{2}^{T}(n) \right\}
\end{bmatrix}_{2M \times 2M}: \text{covariance matrix}
\]

\[
\underline{\mathrm{u}} =
\begin{bmatrix}
\underline{\mathrm{g}}_{2} \\
-\underline{\mathrm{g}}_{1}
\end{bmatrix}_{2M \times 1}: \text{channels vector}
\]

We readily find that:

\[
\left[ \mathrm{R} \right] \underline{\mathrm{u}} = \underline{0}_{2M}
\]

The term $\underline{\mathrm{u}}$ is the eigenvector associated to the zero-valued eigenvalue of $\mathrm{R}$.
If the two impulse responses have no common zeros, there is a \textbf{single zero-valued eigenvalue}.

The last equation is satisfied only for the true covariance matrix.
We can only rely on an \textbf{estimate} of $\left[\mathrm{R}\right]$, thus we search for the \textbf{minimum eigenvalue}.
The eigenvalue decomposition has complexity $\mathcal{O}\left( \left( 2M \right)^{3} \right)$ and considering that $M$ can be large (i.e. channel responses length of thousands of taps), it is convenient to perform an \textbf{iterative search}.

The idea is to \textbf{minimize the norm} of the vector $\left[\mathrm{R}\right] \underline{\mathrm{u}}$.  
We pose a constraint on the norm of $\underline{\mathrm{u}}$ to prevent the trivial solution $\underline{\mathrm{u}} = \underline{0}_{2M}$,  
then solve the constrained minimization problem:

\[
\underline{\hat{\mathrm{u}}} = \underset{\underline{\mathrm{u}}}{\mathrm{arg \ min}} \left( \underline{\mathrm{u}}^{T}  \left[\mathrm{R}\right] \underline{\mathrm{u}}\right),
\qquad
\underline{\mathrm{u}}^{T} \underline{\mathrm{u}} = 1
\]

We can solve the problem using an LMS-like procedure.
We define the error signal:

\[
e\left( n \right) = \frac{\underline{\mathrm{u}}^{T}\left( n \right) \  \underline{\mathrm{y}}\left( n \right)}{\norm{\underline{\mathrm{u}} \left( n \right)}},
\qquad
\underline{\mathrm{y}}\left( n \right) = 
\begin{bmatrix}
\underline{\mathrm{y}}_{1}^{T} \\
\underline{\mathrm{y}}_{2}^{T}
\end{bmatrix}
\]

It can be observed that the \textbf{MSE minimization} of $e\left( n \right)$ corresponds to the constrained minimization problem.
The constrained LMS algorithm is:

\[
\underline{\mathrm{u}}\left( n + 1 \right) = \underline{\mathrm{u}}\left( n \right) - \mu \ e\left( n \right) \nabla e\left( n \right),
\qquad
\nabla e\left( n \right) = \frac{1}{\norm{\underline{\mathrm{u}} \left( n \right)}}
\left[ 
\underline{\mathrm{y}}\left( n \right) - e\left( n \right) \frac{\underline{\mathrm{u}}\left( n \right)}{\norm{\underline{\mathrm{u}} \left( n \right)}}
\right]
\]

After convergence, it can be proved that:

\[
\left[ \mathrm{R} \right] \frac{\underline{\mathrm{u}}\left( \infty \right)}{\norm{\underline{\mathrm{u}}\left( \infty \right)}} =
E \left\{ e^{2}\left( \infty \right) \right\} \frac{\underline{\mathrm{u}}\left( \infty \right)}{\norm{\underline{\mathrm{u}}\left( \infty \right)}}
\]

This confirms that the iterative algorithm converges to the correct solution:
$E \left\{ e^{2}\left( \infty \right) \right\} $ is the smallest eigenvalue of $\left[ \mathrm{R} \right]$.

We are now ready for TDOA estimation:

\begin{enumerate}
    \item extract the estimated channel responses $\hat{g}_{1}(n)$ and $\hat{g}_{2}(n)$ from $\underline{\mathrm{u}}\left( \infty \right)$

    \item estimate the propagation delays as:

    \[
    \hat{\tau}_{k} = \underset{\tau}{\mathrm{arg \ max}} \left( \hat{g}_{k}\left( \tau \right) \right),
    \qquad
    k = 1, 2
    \]

    \item estimate the TDOA as:

    \[
    \hat{\tau}_{12} = \hat{\tau}_{2} - \hat{\tau}_{1}
    \]
\end{enumerate}

\clearpage
